%  LaTeX support: latex@mdpi.com 
%  For support, please attach all files needed for compiling as well as the log file, and specify your operating system, LaTeX version, and LaTeX editor.

%=================================================================
\documentclass[climate,article,submit,pdftex,moreauthors]{Definitions/mdpi} 

%=================================================================
%----------
% journal

% MDPI internal commands
\firstpage{1} 
\makeatletter 
\setcounter{page}{\@firstpage} 
\makeatother
\pubvolume{1}
\issuenum{1}
\articlenumber{0}
\pubyear{2025}
\copyrightyear{2025}
%\externaleditor{Academic Editor: Firstname Lastname}
\datereceived{25 February 2025} 
%\daterevised{} % Only for the journal Acoustics
\dateaccepted{} 
\datepublished{} 
\hreflink{https://doi.org/} % If needed use \linebreak
%\doinum{}

%=================================================================
% Add packages and commands here. The following packages are loaded in our class file: fontenc, inputenc, calc, indentfirst, fancyhdr, graphicx, epstopdf, lastpage, ifthen, lineno, float, amsmath, setspace, enumitem, mathpazo, booktabs, titlesec, etoolbox, tabto, xcolor, soul, multirow, microtype, tikz, totcount, changepage, attrib, upgreek, cleveref, amsthm, hyphenat, natbib, hyperref, footmisc, url, geometry, newfloat, caption

\graphicspath{{figures/}} % path to folder with figures
\usepackage{subcaption}
\usepackage{listings}
\usepackage{xcolor}
\usepackage{gensymb} % for \textdegree
\usepackage{tabularx}
\newcolumntype{Y}{>{\centering\arraybackslash}X}
\usepackage{amsmath,mathtools,amsthm}
	\setcounter{MaxMatrixCols}{15}
\usepackage{array}
\usepackage{amssymb}
\usepackage{amsfonts}
\usepackage{float}
\usepackage{chngcntr}
\counterwithin*{equation}{section}
\usepackage[super]{nth}
\usepackage{longtable}
\usepackage{tabularx}
\usepackage{multirow}
\usepackage{adjustbox}

\definecolor{deepblue}{rgb}{0,0,0.5}
\definecolor{deepred}{rgb}{0.6,0,0}
\definecolor{deepmagenta}{RGB}{195,12,214}
\definecolor{deepgreen}{RGB}{29,122,30}
\definecolor{mintgreen}{RGB}{192,249,192}
\definecolor{codepurple}{RGB}{204, 0, 153}
\definecolor{LightCyan}{rgb}{0.88,1,1}
\definecolor{GainsboroPurple}{RGB}{247,176,247}
\definecolor{LightSlateGrey}{RGB}{124,118,118}
%    
\lstdefinestyle{mystyle}{
    commentstyle=\color{LightSlateGrey},
    keywordstyle=\color{deepgreen}\bfseries,
    emphstyle=\color{deepred},
    numberstyle=\tiny\color{deepblue},
    stringstyle=\color{codepurple},
    basicstyle=\ttfamily\scriptsize,
    breakatwhitespace=false,         
    breaklines=true,                 
    captionpos=b,                    
    keepspaces=true,                 
    numbers=left,                    
    numbersep=5pt,                  
    showspaces=false,                
    showstringspaces=false,
    showtabs=false,                  
    tabsize=2
}
\lstset{style=mystyle}


%=================================================================
% Full title of the paper (Capitalized)
\Title{Ensemble Learning Methods of Image Classification Using GRASS GIS for Mapping Seasonal Fluctuations in Etosha Pan, Namibia}

% MDPI internal command: Title for citation in the left column
\TitleCitation{Ensemble Learning Methods of Image Classification Using GRASS GIS for Mapping Seasonal Fluctuations in Etosha Pan, Namibia}

% Author Orchid ID: enter ID or remove command
\newcommand{\orcidauthorA}{0000-0002-5759-1089} % Polina Add \orcidA{} behind the author's name

% Authors, for the paper (add full first names)
\Author{Polina Lemenkova $^{1}$*\orcidA{}}

%\longauthorlist{yes}

% MDPI internal command: Authors, for metadata in PDF
\AuthorNames{Polina Lemenkova}

% MDPI internal command: Authors, for citation in the left column
\AuthorCitation{Lemenkova, P.}
% If this is a Chicago style journal: Lastname, Firstname, Firstname Lastname, and Firstname Lastname.

% Affiliations / Addresses (Add [1] after \address if there is only one affiliation.)
\address{%
$^{1}$ \quad Alma Mater Studiorum -- Università di Bologna, Department of Biological, Geological and Environmental Sciences, Via Irnerio 42, Bologna, Emilia-Romagna, Italy, IT-40126. ORCID- 0000-0002-5759-1089, (P.L.).

$^{2}$ \quad Freie Universität Bozen | Libera Università di Bolzano, Faculty of Agricultural, Environmental and Food Sciences, Universitätsplatz 5, IT-39100, Bolzano, Trentino-Alto Adige | S\"udtirol, Italy.
}

% Contact information of the corresponding author
\corres{Correspondence: polina.lemenkova2@unibo.it ; polina.lemenkova@unibz.it . Tel.: +39-3446928732 (P.L.)}

% Current address and/or shared authorship
%\firstnote{Current address: Affiliation 3.} 

% Abstract (Do not insert blank lines, i.e. \\) 
\abstract{Ephemeric desert lakes exhibit high fluctuation in salinity and water level that depend on the oscillations of climatic parameters – temperature and precipitation. Understanding these rhythms in seasonal lakes of southern Africa is of high interest to a variety of audiences including environmental planners, climate analysts, and policy makers in protected areas. In this pa- per, we present a Machine Learning (ML) application for image analysis for visualization and processing of Earth observation data collected from the United States Geological Survey (USGS). The case study is focused on Etosha pan which is a large salt pan located in northern Namibia, Africa. It is a vast basin in the land surface where water collects seasonally depending on rains and evaporates during dry periods with the surface covered by salt and minerals. Satellite image classification presents an effective method for monitoring these changes using a comparison of several scenes collected during dry and wet periods. While traditional classification methods are usually based on the Geographic Information Systems (GIS) tools, a novel tool of Machine Learning (ML) presents a more effective, automated and accurate approach to image processing. This paper presents the application of one of the ML methods using Random Forest (RF) algorithms of Geographic Resources Analysis Sup- port System (GRASS) GIS which utilises the embedded libraries of Python Scikit-Learn for image processing. The results show landscape dynamics in Etosha Pan from March to August in 2022 using a series of Landsat satellite images classified by RF method.}

% Keywords
\keyword{keyword 1; keyword 2; keyword 3; keyword 4; keyword 5; keyword 6; keyword 7; keyword 8; keyword 9; keyword 10} 

% The fields PACS, MSC, and JEL may be left empty or commented out if not applicable
%\PACS{J0101}
%\MSC{}
%\JEL{}

\begin{document}

\section{Introduction}

The use of Machine Learning (ML) has been in the centre of attention recently for the challenge of Earth observation data processing. ML has become available approach in cartography along with the advances in programming and these methods are increasingly used in various tasks of environmental monitoring and mapping. The use of these tools ensures automation of satellite data processing, supports operative monitoring using time series and has an important effect on accuracy and effectiveness of data processing for cartographic visualization. Thus, high level of automation in data processing achieved by ML approaches makes it possible to map land cover types rapidly using computer vision approaches applied to satellite images.

With numerous applications in the environmental and Earth science sector, many studies are currently investigating the applicability of programming algorithms in processing Remote Sensing (RS) data \cite{10080513,10334172,Lemenkova202227,9964623,jimaging9050098,9393066}. On the one hand, the availability of the RS data supported by the increasingly rising pool of Earth observation sensors and products of satellite missions, both commercial and open, brings new opportunities to the thematic mapping. On the other hand, such challenges require the development of the advanced methods for processing these datasets, as discussed earlier \cite{JASIEWICZ20111525}. In this regards, ML presents significant benefits and opportunities compared to the traditional cartographic methods possible using Geographic Information Systems (GIS). For instance, due to the introduced training and randomness in raster data processing ensured by computer-vision algorithms of automated image analysis, ML enables to avoid subjectivity in image classification. 

This research aims to build a workflow for satellite image processing using ensemble learning techniques of Random Forest (RF) algorithm. Based on integrating several ML approaches, it includes the decision trees algorithm with controlled variance that improves the accuracy of image partition. Technically, the research is performed using the ML-based modules of the Geographic Resources Analysis Support System (GRASS) GIS software using existing methods \cite{lemenkova2025automation}. The aim of this study is to dissect the effects of seasonal water/salt fluctuations in the Etosha Pan of northern Namibia on spatial heterogeneity and patterns of the land cover patches in the surrounding lacustrine landscapes. The RF approach was selected as valuable methods since it uses the computer vision approaches that have the anti-noise capacity which is useful for image processing used for mapping changing in basin from evaporated to flooding period. Such benefits are useful when processing satellite data with high spectral heterogeneity and landscape fragmentation such as northern Namibia. Furthermore, the RF techniques consider spectral properties of data obtained from diverse sensors such as surface forms, texture of land parcels, topology and pattern heterogeneity. This enables to process and analysis complex landscape patterns in diversified environmental regions. In view of this, the use of the ML by GRASS GIS is considered a key step in the development of cartographic tools of the RS data processing.

Ensemble learning approach of image processing is the method of image processing that applies several learning paradigms to achieve high-performance satellite image processing. The effects of the ensemble learning resulting in the improved cartographic outputs approach can be explained by the integrated power of different tools which bring their benefits to the workflow compared to these algorithms when used alone. In ML-based satellite image analysis, the ensemble learning includes a series of alternative models, such as RF or decision trees, ensuring a more adjusted workflow. Specifically in computer vision, using ensemble learning methods of ML is an essential part of automated image analysis and mapping based on these images. Such benefits of ML for cartographic workflow have been reported in some publications in recent decades \cite{YOO2019155,SINGH2021100624,MULIK2023106697}. However, the suitability of specific software that includes these methods as embedded algorithms for image processing requires more attention. In RS techniques, two major methods have been discussed and compared recently, including supervised \cite{GOLDBLATT2018253,MARTINEZMONTOYA2010978,INZANA200359} or unsupervised \cite{SHIH199497,LU2023116898,Lemenkova202322} approaches to satellite image classification.

The state-of-the-art methods of image classification in traditional GIS methods have certain limitations which can be summarised as follows. Firstly, since image classification is generally based on spectral reflectance of pixels, some cases may take place. This is cased by the ground properties of pixels which are not representing real land surface objects and strongly depend on the resolution of satellite images. As a result, objects which have very similar spectral properties can be erroneously assigned to the same land cover class. Secondly, object topology is often avoided in when interpreting the imagery. To avoid this drawback, some studies analyse spatial relationships between the objects using Object Based Image Analysis (OBIA) techniques \cite{LUCIANO2019127,PHIRI2018170,DEWALI20233109,MAO202011}. Nevertheless, this is a different approach that requires specific software and tools. Thirdly, the inherited properties of the land surface such as texture, curvature of patches, context of disposition of several objects representing the scene, shape, form and structure of objects is difficult to properly identified using the traditional image classification techniques in standard GIS \cite{OUSMANOU2024100239,ZHANG2023129590,BAKR2010592}. Moreover, increased variability of images with different properties and resolution require more advanced methods compared to the traditional approaches of image analysis that use pixel-based classification algorithms. This lead to the issues in accuracy and precision of image classification process \cite{Fisher,McKee}.

In view of these limitations, the use of the novel ML methods for RS data processing and programming algorithms for classification approaches has recently drawn much attention \cite{Ayhan}. Among them, Python algorithms of Scikit-Learn library which presents a powerful instrument for ML methods have been investigated in a number of works \cite{9884244,9883588,9006205,7724471}. The RF algorithm embedded in the mapping of crop fields using Landsat imagery was thoroughly described by \cite{LUCIANO2021106063}. \cite{AHMED201589} proposed an integrated application of the RF algorithm based on the multi-source data combining LiDAR and Landsat where connected components in the RS data are particularly well suited for complex forest canopy analysis using data on cover and height. Diverse examples of the RF algorithm applications in various cases of RS data processing include crop mapping for food and agricultural research \cite{OLIPHANT2019110,TELUGUNTLA2018325}, hazard risk assessment using fire mapping \cite{COLLINS2020111839}, urban sprawl analysis \cite{GHOSH201431} and environmental monitoring of land cover types \cite{MELVILLE201846,GRINAND201368}. 

More research projects have demonstrated the potential of the RF approach of ML to image processing \cite{COLLINS2020111839,WEI2020112005,WANG202294}. Nevertheless, it still suffers from some limitations regarding its application in the GRASS GIS. For example, RF algorithm needs to be better exploited using GRASS GIS scripting modules in view of the current lack of a practical framework and hands-on applications in case studies. There have been limited case studies on the use of GRASS GIS using scripts for image processing in environmental applications \cite{STRIGARO201668,geomatics5010005,ROCCHINI2017166,7326254}. In further examples, \cite{7326254} evaluated the use of GRASS GIS modules for marine geological mapping using MODIS/Aqua RS data in coastal areas. The approach presented by \cite{6407912} reports the links between GRASS GIS and Python and performs relatively well for environment monitoring using satellite imagery but the proper scripts of GRASS GIS are not reported since the major point was using Python for launching GRASS GIS. Moreover, few experimental studies have thoroughly evaluated the GRASS GIS modules for diverse ecological applications and mapping. However, the method of using RF by the GRASS GIS is not reported in the existing available works and needs evaluation. 

The selection of methods for image processing and software functionality which enables to process the RS data are the most important factors for studies using Earth observation data analysis. In this research, we present a RF-based classification method using GRASS GIS software integrated with Python library Scikit-Learn, which uses  of ML algorithms as a background tools for data processing. Unlike existing methods of image classification, the ensemble learning method of RF considers the diversity among the spectral reflectance of pixels when modelling topology between landscape patches. Other advantages of this method is that ensemble learning uses data on texture and spectral features to measure the similarity of neighbouring landscape patches during the process of generating random decision trees. 

The goal of this study is to apply the scripting techniques of GRASS GIS with the objective to assess the novel method of ensemble learning with a case of RF algorithms applied to cartographic workflow. In this way, we evaluate its functionality of the machine learning tools of GRASS GIS and their applicability in RS data processing to classification of multispectral satellite images. A motivation of this study is to perform environmental analysis of the Etosha salt pan, northern Namibia, with the aim to explore the climatic and seasonal factors that affect the water availability in the basin. The paper is organised as follows. Chapter 2 explains the ensemble learning method with RF approach to RS data processing. The methodology is presented in snippets of GRASS GIS code. In the following section, the results of image classification are presented. Finally, the last section concludes the paper with the discussion regarding the landscapes and climate-related land cover changes in the Etosha Pan, northern Namibia.

\section{Study Area}
The study area is focused on the saline Etosha Pan, Figure \ref{fig01}. 
 
 \setlength{\belowcaptionskip}{-10pt}
 \begin{figure}[H]
	\noindent\includegraphics[width=\textwidth]{Fig_01.jpg}
\caption{Topographic map of Namibia with indicated study area: Etosha Pan. Mapping software: Generic Mapping Tools (GMT). Map source: author.}
\label{fig01}
\end{figure}

Etosha Pan belongs to the Cuvelai-Etosha Basin, a transboundary wetland area within the topographic depression located between the southern Angola and northern Namibia. The basin developed from the Palaeolake Etosha since Tertiary and presents the depression that collects waters from three drainage systems off the Angolan highlands, the Cubango and the Cuvelai Drainage System \cite{MILLER}. The topography of the region is generally flat which causes the diurnal oscillation of wind and variations in thermo-topographical gradients \cite{Preston}. Highly changing hydro-geomorphological setting of the deltaic drainage system of Etosha Pan lead to the regular fluctuations of the water level where lake changes from a basing flooded by water to a completely dry pan covered by crust salt and minerals \cite{Hipondoka}. 

Climate-related hazards such as frequent and severe droughts, unstable water inflow and extreme temperatures typical for arid and semi-arid climate of Namibia affect rare species that use Etosha as shelter habitat \cite{cli5030051}. To preserve the world-renowned resources of vulnerable and fragile ecosystem around Etosha Pan, nature conservation measures and policies in Namibia were undertaken. This lead to the creation of the natural reserve of the Etosha National Park aimed at regulating the negative effects from agriculture activities, livestock husbandry, overgrazing, and tourism \cite{Gargallo,Novelli}. Landscapes of the Etosha National Park are dominated by the distribution of dry savannah which is a typical environmental characteristics of Namibia. Nevertheless, the variations of vegetation classes is highly fragmented which depends on the inter‐seasonal variations, availability of moisture and correlate with the actual rainfall situation \cite{Wagenseil}. The detailed land cover types are visualised on the map in Figure \ref{fig02}.

 \begin{figure}[H]
	\noindent\includegraphics[width=\textwidth]{Fig_02.jpg}
\caption{Land cover types of Namibia according to  Food and Agriculture Organizatio (FAO). Mapping software: QGIS. Map source: author.}\setlength{\belowcaptionskip}{-30pt}
\label{fig02}
\end{figure}

Oscillations in hydrological patterns from flooding to complete desiccation indicate the progressively increasing aridity and desertification of the Etosha basin affected by regional droughts, impacts from the Kalahari Desert and climate change. Two distinct seasons - dry from April to September and wet from October to March - distinctly divide the climate and environmental patterns while the intra-seasonal water influx is unstable and depend on the ephemeral streams forming occasional pools \cite{Himmelsbach}. At the same time, unique landscapes of Etosha create a habitat for ecosystems used by migrating birds and rare species as habitats \cite{Fox,LeRoux,Berry}. The ability of wetland birds such as flamingos to predict rain fronts and showers using changes in pressure gradients enables them to rapidly find water available in ephemeral reservoirs and pools of Etosha. This results in the significant number of migrating rare birds that use the Etosha as a destination place during severe droughts in other regions \cite{Simmons}. 

%----------- section ----------->
\section{Materials and Methods}

%----------- subsection ----------->
\subsection{Data}

\begin{figure}[H]
  \centering
  \begin{tabular}[b]{c}
    \includegraphics[width=2.9cm]{dummy.jpg} \\
    \small (a): 
  \end{tabular}
  \begin{tabular}[b]{c}
    \includegraphics[width=2.9cm]{dummy.jpg} \\
    \small (b): 
  \end{tabular}
    \begin{tabular}[b]{c}
    \includegraphics[width=2.9cm]{dummy.jpg} \\
    \small (c): 
  \end{tabular}
  \begin{tabular}[b]{c}
    \includegraphics[width=2.9cm]{dummy.jpg} \\
    \small (d):
  \end{tabular}
  \caption{Original data: Landsat 8-9 OLI/TIRS images on XXX collected on years XXX, XXX, XXX and XXX. Source: USGS \cite{Landsat}. Compilation source: author.}
  \label{fig04}
\end{figure}

\section{Materials and Methods}
 
In the workflow of image processing, the Landsat images are computed by a GRASS GIS software which presents a powerful toolset for image analysis \cite{GRASS_GIS_software}. Its excellent and diversified functionality in image processing are reflected in many reported studies which use Earth observation data processed by diverse modules of GRASS GIS \cite{HOFIERKA2009387,lemenkova2024exploitation,Lemenkova202320,ROCCHINI201382,Lemenkova202313,JASIEWICZ20111162}. The Generic Mapping Tools (GMT) \cite{Wessel} is adopted to map topographic information and visualise the location of the study area using overlaid extent of the multispectral Landsat images, Figure \ref{fig01}. More specifically, the GMT was used using scripting techniques reported and described in earlier works \cite{Lemenkova202203,Lemenkova202217}. Here the extent of the images is denoted by the rotated square showing the location of the Landsat scenes. Their coordinates were derived from the metadata and correspond to the technical parameters of the Landsat mission. The QGIS software was used for land cover mapping to represent the geospatial information of the landscapes in northern Namibia, Figure \ref{fig02}. 

\subsection{Data}
The original images are visualised in Figure \ref{fig03}. Major technical characteristics that differ for each Landsat scene are provided in Table \ref{tab01}. The data were obtained from the EarthExplorer repository of the United States Geological Survey (USGS) and included six satellite images of the Landsat 8-9 Operational Land Imager and Thermal Infrared Sensor (Landsat 8-9 OLI/TIRS). 

\setlength{\belowcaptionskip}{-30pt}
\begin{figure}[H]
  \centering
  \begin{tabular}[b]{c}
    \includegraphics[width=4cm]{Fig_03a_mar.jpg} \\
    \small (a)
  \end{tabular}
  \begin{tabular}[b]{c}
    \includegraphics[width=4cm]{Fig_03b_apr.jpg} \\
    \small (b)
  \end{tabular}
  \begin{tabular}[b]{c}
    \includegraphics[width=4cm]{Fig_03c_may.jpg} \\
    \small (c)
     \end{tabular}
     \begin{tabular}[b]{c}
    \includegraphics[width=4cm]{Fig_03d_jun.jpg} \\
    \small (d)
  \end{tabular}
  \begin{tabular}[b]{c}
    \includegraphics[width=4cm]{Fig_03e_jul.jpg} \\
    \small (e)
  \end{tabular}
  \begin{tabular}[b]{c}
    \includegraphics[width=4cm]{Fig_03f_aug.jpg} \\
    \small (f)
  \end{tabular}
  \caption{Landsat 8-9 OLI/TIRS images used as dataset covering the Etosha Pan region, Namibia, 2022: \textbf{(a)}: 26 March;\textbf{(b)}: 11 April; \textbf{(c)}: 5 May; \textbf{(d)}: 6 June; \textbf{(e)}: 16 July; \textbf{(f)}: 9 August. The images are shown in RGB colours. Data source: USGS.}
  \label{fig03}
\end{figure}
\vspace{2em}

The colour composites of the dataset are presented in Figure \ref{fig04} with different combinations used as the RGB triplets from the Landsat bands. The raw dataset illustrates the environmental changes in Etosha Pan that fluctuates from the wet period when the lake is filled with water to the dry period when the basin becomes covered with a salt crust. 


%----------- subsection ----------->
\subsection{Subsection}

%----------- subsection ----------->
\subsection{Subsection}

%----------- section ----------->
\section{Results}

\begin{figure}[H]
  \centering
  \begin{tabular}[b]{c}
    \includegraphics[width=6.3cm]{dummy.jpg} \\
    \small (a)
  \end{tabular}
  \begin{tabular}[b]{c}
    \includegraphics[width=6.3cm]{dummy.jpg} \\
    \small (b) 
  \end{tabular}
    \begin{tabular}[b]{c}
    \includegraphics[width=6.3cm]{dummy.jpg} \\
    \small (c)
  \end{tabular}
  \begin{tabular}[b]{c}
    \includegraphics[width=6.3cm]{dummy.jpg} \\
    \small (d) 
  \end{tabular}
  \caption{Results. Image processing source: author.}
  \label{fig06}
\end{figure}


%----------- section ----------->
\section{Discussion}

%----------- section ----------->
\section{Conclusions}

%----------- section ----------->
\authorcontributions{Supervision, conceptualization, methodology, software, resources, funding acquisition, and project administration, writing---original draft preparation, methodology, software, data curation, visualization, formal analysis, validation, writing---review and editing, and investigation, P.L.}

\funding{The publication was funded by the Editorial Office of XXXX, Multidisciplinary Digital Publishing Institute (MDPI), by providing 100\% discount for the APC of this manuscript.}

\institutionalreview{Not applicable.}

\informedconsent{Not applicable.}

\dataavailability{Not applicable} 

\acknowledgments{The authors thank the reviewers for reading and review of this manuscript.}

\conflictsofinterest{The authors declare no conflict of interest.} 

\abbreviations{Abbreviations}{
The following abbreviations are used in this manuscript:\\

\noindent 
\begin{tabular}{@{}ll}
MDPI & Multidisciplinary Digital Publishing Institute\\
DOAJ & Directory of open access journals\\
LD & Linear dichroism
\end{tabular}
}

%%%%%%%%%%%%%%%%%%%%%%%%%%%%%%%%%%%%%%%%%%
%% Optional
\appendixtitles{no} % Leave argument "no" if all appendix headings stay EMPTY (then no dot is printed after "Appendix A"). If the appendix sections contain a heading then change the argument to "yes".
\appendixstart
\appendix
\section[\appendixname~\thesection]{}
\subsection[\appendixname~\thesubsection]{}
The appendix 

\section[\appendixname~\thesection]{}
All appendix sections must be cited in the main text. In the appendices, Figures, Tables, etc. should be labeled, starting with ``A''---e.g., Figure A1, Figure A2, etc.

%%%%%%%%%%%%%%%%%%%%%%%%%%%%%%%%%%%%%%%%%%
\begin{adjustwidth}{-\extralength}{0cm}
%\printendnotes[custom] % Un-comment to print a list of endnotes

\reftitle{References}
%\nocite{*}

%=====================================
% References, variant A: external bibliography
%=====================================
\bibliography{Bibliography.bib}

%%%%%%%%%%%%%%%%%%%%%%%%%%%%%%%%%%%%%%%%%%
\end{adjustwidth}
\end{document}
